\subsection*{When a review loop should stop}

The third question was when a revision loop should stop, given that reviewers never run out of concerns. We recorded 24 real loops (12 flawed items $\times$ 2 configurations), each with five revisions and six panel reviews, and judged the state of every planted flaw after every revision (Fig~\ref{fig4}).

\textbf{Planted flaws were fixed early; new errors accumulated later.} The first revision resolved most planted flaws: all four were resolved in 83\% of items with gpt-5.6-sol and 67\% with gpt-5.6-luna, leaving 0.17 and 0.33 unresolved flaws per item. Later revisions added little (0.08 unresolved flaws per item after four revisions in both configurations), and no planted flaw was resolved by asserting an analysis that had not been done. Every flaw that remained unresolved after a revision was a false statement about the model or its tokenization, the flaw type that reviewers also rarely detected in the benchmark; the panel had not flagged these statements, which most likely explains why the executor left them unchanged. Meanwhile the judge found newly introduced errors, and these accumulated: from 0.17 and 0.67 per item after the first revision to 1.75 and 1.83 after the fourth, the platform's default budget, and 1.00 and 1.92 after the fifth (Fig~\ref{fig4}B). With 12 items per configuration we do not read the fall for gpt-5.6-sol at the fifth revision as a reversal; in both configurations the write-ups still contained far more new errors than after the first revision. The two configurations erred differently. Statements that presented an analysis, control or result as done although the original write-up did not contain it (labelled fabrications by the judge) made up about half of the new errors with gpt-5.6-sol from the third revision onwards (0.75 of 1.25 per item after the third revision and 0.92 of 1.75 after the fourth), although the executor had been instructed not to report new results; with gpt-5.6-luna they stayed at or below 0.17 per item. Most of these statements misdescribed what the write-up had done, for example by crediting it with procedures or thresholds it never contained, rather than reporting new numerical results. Write-ups grew from about 6,200 characters to 25,000--56,000 after the fourth revision, mostly through added caveats and plans for pending analyses.

\textbf{The panel's signals separate the original from the revised write-ups but do not track later changes.} The number of critical and high critiques fell by about half after the first revision (from 23.8 to 11.8 per review with gpt-5.6-sol and from 16.3 to 8.8 with gpt-5.6-luna). Pooled over all six reviews of each loop, it therefore correlated with the number of unresolved planted flaws (Spearman $\rho$ 0.54 with gpt-5.6-sol and 0.49 with gpt-5.6-luna; 0.66 and 0.70, respectively, for critical critiques alone). This correlation reflects the contrast between the original write-up, in which all four flaws were unresolved, and the revisions; after the first revision the count neither followed the few remaining planted flaws nor rose as new errors accumulated, but levelled off at 7--12 per review and never reached zero (Fig~\ref{fig4}C). The merged grade was F for 23 of the 24 original write-ups, rose to D for all 24 first revisions and was D in 113 of the 120 reviews of revised write-ups (6 were F and 1 was C). These grade levels depend on the merge: the loops' panels merged at $\tau=0.5$, which escalated few agreements, whereas with effective merging the benchmark panels graded every item F (see Merging paraphrased critiques). Because critical and high critiques never disappeared and consecutive executor outputs never became similar enough, every candidate rule that requires the absence of high-severity critiques or a stable executor output never stopped within five revisions, with or without the advancement gate. The only adaptive signal that fired was grade stability, which in practice detects that the grade has stopped changing: it stopped after a mean of 3.0 revisions and stopped prematurely---while a planted flaw was still unresolved---in 3 of 24 runs (Fig~\ref{fig4}D). A fixed budget of three revisions performed identically (3 of 24), and a fixed budget of four revisions stopped prematurely in 2 of 24 runs. Applied to the 24 runs together, the pre-specified selection rule therefore sets the platform default to a budget of four revisions after the initial submission, followed by a final review; the adaptive rule remains available. The budget rules use only the revision count and therefore do not depend on the merge method. Within configurations the choice is less clear-cut: with gpt-5.6-luna, four revisions (1 of 12 premature stops) did better than three revisions and grade stability (2 of 12 each), whereas with gpt-5.6-sol both budgets and grade stability each stopped prematurely in 1 of 12 runs, and the tie-break alone would select three revisions. The difference between the candidates rests on one run of 24 and comes entirely from flaws the panel does not detect, so no stopping signal derived from the panel could have prevented those premature stops. The criterion also does not count introduced errors, and by that measure four revisions is the most costly eligible rule: 1.75 and 1.83 new errors per item, against 1.25 and 1.58 after three revisions and 0.17 and 0.67 after one.

\textbf{Interpretation.} With current LLM reviewers, stopping is a budget decision rather than a convergence decision: the panel's assessment improves sharply after the first revision and then plateaus at a level that would never pass an absolute threshold, while additional revisions trade a small reduction in unresolved planted flaws for an accumulation of new errors that the panel's critique counts did not register. The letter grade can show this trajectory only when merging is weak; with effective merging it saturates at F, so it supports neither use as an acceptance criterion nor use as a stopping signal. A loop's output should accordingly be treated as a draft for human acceptance. In these loops, in which the executor could only revise text, revisions beyond the first removed fewer planted flaws (0.08--0.25 per item between the first revision and the fourth) than the new errors they introduced (1.2--1.6 per item). The default of four revisions therefore minimises planted flaws left unresolved at the price of more introduced errors. When introduced errors weigh more, a budget of one revision is preferable: it left 0.17 and 0.33 unresolved planted flaws and introduced 0.17 and 0.67 new errors per item. Executed loops may behave differently, because a revision can rerun a corrected analysis rather than only reword a claim; in the executed case study, confirmed defects were still being fixed in the fourth revision (see Executed agent runs). For these reasons the budget is a run-level setting.
